\documentclass[techrep]{ipsj}
%\documentclass{ipsj}
\usepackage{float}
\usepackage[dvipdfmx]{graphicx}
\usepackage{latexsym}
\usepackage{algorithm}
\usepackage{algpseudocode}
\usepackage{amsmath}

\def\Underline{\setbox0\hbox\bgroup\let\\\endUnderline}
\def\endUnderline{\vphantom{y}\egroup\smash{\underline{\box0}}\\}
\def\|{\verb|}

\setcounter{巻数}{59}
\setcounter{号数}{1} 
\setcounter{page}{1}

\受付{2016}{3}{4}
\採録{2016}{8}{1}

\begin{document}

\title{空間熱干渉を考慮した状態推定に基づく\\
動的熱マージン制御手法}
\etitle{A Dynamic Thermal Margin Control Method Based on State Estimation Considering Spatial Thermal Interference}

\affiliate{TMCIT}{東京都立産業技術高等専門学校 \\
Tokyo Metropolitan College of Industrial Technology, Tokyo 140--0011, Japan}

\author{鈴木 黎映}{Reo Suzuki}{TMCIT}
\author{知念 賢一}{Kenichi Chinen}{TMCIT}

\begin{abstract}
エッジコンピューティングの進展に伴い，シングルボードコンピュータを高密度に配置したマイクロクラスタの利用が拡大している．このような環境では，隣接ノードの発熱が自ノードの温度に影響を与える空間熱干渉が発生する．しかし，この外部熱影響は負荷変動などにより時間的に変化するため，自ノードの負荷のみに基づく従来モデルでは予測誤差が増大するという課題がある．
本研究では，外部熱影響を時間的に変動する状態として明示的にモデル化した拡張熱モデルを提案する．さらにカルマンフィルタを適用することで，自ノード温度と外部熱影響を同時に推定し，空間熱干渉を考慮した温度予測を実現する.
評価実験の結果，提案手法は従来の単一ノードモデルおよび誤差統計に基づく手法と比較して予測精度が向上することを確認した．また，不確実性に基づく動的マージン調整についても初期的な検証を行ったが，現在まだ効果は認められない.
\end{abstract}

\begin{jkeyword}
エッジコンピューティング，熱管理，空間熱干渉，温度予測
\end{jkeyword}

\begin{eabstract}
With the advancement of edge computing, micro-clusters composed of densely deployed single-board computers are increasingly used. In such environments, spatial thermal interference occurs, where heat generated by neighboring nodes affects the temperature of a target node. However, since this external thermal effect varies over time due to workload fluctuations, conventional models based only on local workload suffer from increased prediction errors.
In this study, we propose an extended thermal model that explicitly represents the external thermal effect as a time-varying state. By applying a Kalman filter, both node temperature and external thermal effects are simultaneously estimated, enabling temperature prediction that accounts for spatial thermal interference.
Experimental results show that the proposed method improves prediction accuracy compared with conventional single-node models and error-statistics-based methods. In addition, we conducted a preliminary evaluation of uncertainty-based dynamic margin adjustment; however, no significant effect has been observed at this stage.
\end{eabstract}

\begin{ekeyword}
edge computing, thermal management, spatial thermal interference, temperature prediction
\end{ekeyword}
\maketitle

\section{はじめに}

IoTや5G技術の進展に伴い，データの生成源に近い場所でリアルタイムに演算処理を行うエッジコンピューティングが不可欠となっている\cite{shi2016edge}．
特に，Raspberry Pi 等のシングルボードコンピュータ（SBC）を複数台集積したマイクロクラスタは，
低コストかつ柔軟に構成可能な小規模計算基盤として注目されている\cite{iridispi2014}．
このような環境では，単一ノードの計算能力を補うための分散実行や，障害耐性を高めるためのノード間冗長化が行われる．
また，エッジ側でのリアルタイム処理の重要性が増すにつれて，限られた計算資源を効率的に活用するための制御手法の重要性も高まっている．

一方で，エッジ環境ではサーバ室のような十分な空調設備を前提とできず，限られた空間に複数ノードを集積して運用することが多い．
そのため，筐体内部では排熱の滞留や再循環が発生しやすく，隣接ノードの発熱が自ノード温度に影響する空間熱干渉が問題となる．
この現象は，自ノードのCPU使用率や処理内容のみでは説明できない温度変動を引き起こし，従来の単一ノード熱モデルの予測精度を低下させる要因となる．
さらに，このような外部影響は時間的にも変動するため，単純な定数外乱として扱うことが難しいという特徴を持つ．
また，エッジデバイスは外気温や周囲環境の影響を強く受け，温度変動が性能や消費電力に大きな影響を与えることが報告されている\cite{edgeengine2023}．

従来の熱管理手法では，熱RCモデル等を用いて将来温度を予測し，その予測値に基づいてDVFSやタスク制御を行う．
しかし，こうした手法の多くは自ノード内部の発熱を主対象としており，空間的な熱干渉のような外部影響を明示的に扱っていない．
そのため，隣接ノードが高負荷となる状況では，予測誤差が急増し，固定的な安全マージンでは過小評価にも過大評価にもなりうる．
前者は温度逸脱を招き，後者は不要な性能低下を引き起こす．
このように，外乱を適切に扱えないことが，従来手法の限界となっている．

本研究では，この問題に対して，隣接ノードからの熱影響を単なる未知外乱ではなく，時間的に緩やかに変動する状態量としてモデル化する．
さらに，カルマンフィルタを適用することで，自ノード温度と外部熱影響を同時に推定し，その誤差共分散に基づいて予測の不確実性を定量化する．
そして，不確実性が大きいときには安全マージンを拡大し，小さいときには縮小することで，温度逸脱の抑制とリソース利用効率の両立を目指す．

\section{関連研究}

\subsection{高密度環境における熱干渉}

高密度なコンピューティング環境における熱課題として，ノードから排出された熱風の一部が再び吸気口へと回り込む熱再循環現象がTangら\cite{tang2008sensor}によって指摘されている．
また，データセンタにおける冷却設計や温度管理に関する研究でも，再循環や局所的ホットスポットが温度制御上の重要課題であることが報告されている\cite{moore2005cool}．
これらの研究は，大規模環境における熱問題を対象としているが，小規模なエッジ環境においても同様の現象が発生することが指摘されている．

図\ref{fig:recirculation}に，高密度配置された筐体内における熱流挙動の概念図を示す．
隣接ノードから排出された熱風の一部が筐体内に滞留し，側壁や障害物に沿って吸気側へ回り込むことで，自ノードの吸気温度が上昇する．
この影響は，ノード間距離，筐体構造，ファンの有無，周辺ノードの負荷変動などに依存して時間的に変化するため，一定値の外乱として扱うだけでは不十分である．
さらに，熱の伝播には遅れが存在するため，瞬間的な影響ではなく，時間的に蓄積される形で温度変化として現れる点も重要である．
\begin{figure}[tbp]
    \centering
    \includegraphics[scale=0.37]{cite1.png}
    \caption{高密度配置されたノード間で発生する熱の再循環現象の概念図}
    \label{fig:recirculation}
\end{figure}

\subsection{熱管理における予測制御}

熱管理においては，将来の温度を予測した上で制御を行う予測型アプローチが広く用いられている．
例えば，熱RCモデルに基づく温度推定や，モデル予測制御（MPC）を用いた性能と温度制約の両立が提案されている\cite{zanini2009multicore}．
これらの手法は，温度制約を破らない範囲で性能を最大化できる点で有用である．

しかし，予測制御は予測モデルの妥当性に強く依存する．
モデルが想定していない外乱が加わった場合には，将来温度を過小評価する可能性がある．
特に，本研究で対象とするような高密度マイクロクラスタでは，隣接ノードの発熱が時々刻々と変化するため，自ノード単独の熱モデルのみでは十分な精度を維持することが難しい．

\subsection{不確かさを考慮した動的マージン管理}

予測の不完全さに対処するため，安全余裕を導入する研究も行われている．
また，不確実性を明示的に扱う制御手法として，
確率的モデル予測制御（SMPC）が提案されている\cite{mesbah2016smcp}．
他にも坂本ら\cite{sakamoto2018kbse}は，予測誤差の統計的性質に基づいて不確かさを表現し，それに応じて動的にマージンを変更する考え方を示している．
このような手法は，不要に大きな固定マージンを常に維持する場合と比べて，保守性と効率のバランスを取りやすい．

ただし，既存研究の多くは，予測対象システム内部の時系列変動や負荷変動を主に扱っており，隣接ノードからの熱流入のような空間的かつ物理的な外部干渉を状態推定の対象として組み込んだ例は限定的である．
本研究は，この点に着目し，空間熱干渉を明示的に状態量として扱うことで，不確実性の原因そのものをモデル内で取り込むことを目指す．

\section{提案手法}

本章では，空間熱干渉を考慮した状態推定に基づく動的熱マージン制御手法について述べる．
従来の熱管理手法では，自ノードの負荷のみに基づいて温度を予測し，固定的な安全マージンを設けることで温度逸脱を防ぐことが一般的である．
しかし，高密度に配置されたマイクロクラスタ環境では，隣接ノードの発熱が自ノード温度に影響を与えるため，自ノードのみを対象としたモデルでは予測誤差が増大しやすい．

特に，隣接ノード負荷の変動に起因する熱影響は時間的に変化するため，固定マージンでは過度に保守的になる場合と，不十分になる場合の両方が生じうる．
そこで本研究では，隣接ノードからの熱影響を単なる外乱ではなく，時間的に緩やかに変動する状態量としてモデル化する．
さらに，カルマンフィルタを用いて自ノード温度と外部熱影響を同時に推定し，その推定誤差共分散から予測の不確実性を定量化する．
そして，不確実性に応じて安全温度閾値に対するマージンを動的に調整することで，温度逸脱の抑制とリソース利用効率の両立を図る．

提案手法の全体構成を図\ref{fig:agent}に示す．

\begin{figure}[tbp]
    \centering
    \includegraphics[scale=0.4]{cite2.png}
    \caption{提案手法の構成}
    \label{fig:agent}
\end{figure}

本手法は，温度・負荷観測，負荷予測・状態推定,不確実性の定量化，動的マージン決定の4段階から構成される．
まず，自ノードの温度および負荷，ならびに隣接ノード温度を取得し，状態空間モデルに入力する．
次に，カルマンフィルタにより現在温度と外部熱影響を推定する．
その後，推定結果から将来温度を予測し，予測誤差共分散に基づいて予測の信頼度を評価する．
さらに，その不確実性に応じて安全マージンを決定する．

\subsection{状態空間モデルと状態推定}

自ノード温度は，自身の負荷だけでなく隣接ノードからの熱影響を受ける．
本研究では，この外部熱影響を状態変数として含む拡張熱モデルを構築する．
本モデルは，熱の時間遅れおよび空間的伝播を簡潔に表現可能な線形モデルとして構成している．

状態ベクトルを以下のように定義する．

\begin{equation}
\mathbf{x}_t =
\begin{bmatrix}
T_t \\
H_t
\end{bmatrix}
\end{equation}

ここで，$T_t$ は時刻 $t$ における自ノード温度，$H_t$ は隣接ノードに起因する外部熱影響を表す内部状態である．

また，隣接ノード温度 $z_t^{\mathrm{nbr}}$ から基準温度 $z_{\mathrm{base}}$ を差し引いた偏差

\begin{equation}
z_t = z_t^{\mathrm{nbr}} - z_{\mathrm{base}}
\end{equation}

を外部熱影響の入力として用いる．
すなわち，$z_t$ は隣接ノードの観測温度であり，$H_t$ はその影響を時間的に平滑化した外部熱状態を表す．

システムの状態方程式は以下で与える．

\begin{equation}
\mathbf{x}_{t+1}
=
\mathbf{A}\mathbf{x}_t
+
\mathbf{B}u_t
+
\mathbf{E}z_t
+
\mathbf{d}
+
\mathbf{w}_t
\end{equation}

\begin{equation}
\mathbf{A}
=
\begin{bmatrix}
a & c \\
0 & \rho
\end{bmatrix},
\quad
\mathbf{B}
=
\begin{bmatrix}
b \\
0
\end{bmatrix},
\quad
\mathbf{E}
=
\begin{bmatrix}
0 \\
1-\rho
\end{bmatrix},
\quad
\mathbf{d}
=
\begin{bmatrix}
d \\
0
\end{bmatrix}
\end{equation}

ここで，$u_t$ は自ノードのCPU負荷入力，$\mathbf{w}_t$ はプロセスノイズである．
係数 $a$ は温度の自己回帰的な遷移，$b$ は自ノード負荷の寄与，$c$ は外部熱影響が自ノード温度に与える結合の強さを表す．
また，$\rho$ は外部熱影響の時間的持続性を表す減衰係数であり，$0<\rho<1$ の範囲で設定する．
$\rho$ が大きいほど，外部熱影響は急激には変化せず，過去の状態を保持しやすいことを意味する．

上式を各成分で書くと，次式となる．

\begin{equation}
T_{t+1}
=
aT_t + bu_t + cH_t + d + w_t^{(T)}
\end{equation}

\begin{equation}
H_{t+1}
=
\rho H_t + (1-\rho)z_t + w_t^{(H)}
\end{equation}

すなわち，自ノード温度は自ノードの熱慣性と負荷入力に加え，推定された外部熱影響から決定される．
一方，外部熱影響は隣接ノード温度偏差 $z_t$ を入力としつつ，時間的に平滑に変動する．

観測方程式は以下で表される．

\begin{equation}
y_t
=
\mathbf{C}\mathbf{x}_t + v_t,
\quad
\mathbf{C}
=
\begin{bmatrix}
1 & 0
\end{bmatrix}
\end{equation}

ここで，$y_t$ は観測温度，$v_t$ は観測ノイズである．
観測できるのは温度のみであり，外部熱影響 $H_t$ は直接観測できない．
そこで，本研究ではカルマンフィルタを用いて，温度観測値とモデル予測を逐次統合し，隠れ状態としての外部熱影響を推定する．
カルマンフィルタは状態空間モデルに基づき，観測値と予測値を統合することで最小二乗誤差の意味で最適な状態推定を行う手法である~\cite{kalman1960}．
予測ステップは以下で与えられる．

\begin{equation}
\hat{\mathbf{x}}_{t+1|t}
=
\mathbf{A}\hat{\mathbf{x}}_{t|t}
+
\mathbf{B}u_t
+
\mathbf{E}z_t
+
\mathbf{d}
\end{equation}

\begin{equation}
\mathbf{P}_{t+1|t}
=
\mathbf{A}\mathbf{P}_{t|t}\mathbf{A}^{\top}
+
\mathbf{Q}_t
\end{equation}

更新ステップは以下で与えられる．

\begin{equation}
\mathbf{K}_{t+1}
=
\mathbf{P}_{t+1|t}\mathbf{C}^{\top}
\left(
\mathbf{C}\mathbf{P}_{t+1|t}\mathbf{C}^{\top}
+
\mathbf{R}
\right)^{-1}
\end{equation}

\begin{equation}
\hat{\mathbf{x}}_{t+1|t+1}
=
\hat{\mathbf{x}}_{t+1|t}
+
\mathbf{K}_{t+1}
\left(
y_{t+1}
-
\mathbf{C}\hat{\mathbf{x}}_{t+1|t}
\right)
\end{equation}

\begin{equation}
\mathbf{P}_{t+1|t+1}
=
\left(
\mathbf{I}
-
\mathbf{K}_{t+1}\mathbf{C}
\right)
\mathbf{P}_{t+1|t}
\end{equation}

ここで，$\mathbf{Q}_t$ はプロセスノイズ共分散，$\mathbf{R}$ は観測ノイズ共分散である．
本研究の拡張版では，隣接ノード温度の変化が大きい区間では外部熱影響の不確かさも増大すると考え，$H_t$ に対応するプロセスノイズを時変とする．
具体的には，

\begin{equation}
q_H(t)
=
q_H^{(0)} + \beta_{qh} |z_t - z_{t-1}|
\end{equation}


\begin{equation}
\mathbf{Q}_t
=
\begin{bmatrix}
q_T & 0 \\
0 & q_H(t)
\end{bmatrix}
\end{equation}

を用いる．
ここで，$q_T$ は温度状態に対応するプロセスノイズ，$q_H^{(0)}$ は外部熱影響状態の基礎ノイズ，$\beta_{qh}$ は隣接温度変化に対する感度係数である．
これにより，隣接ノード温度が急激に変化する区間では，共分散が大きくなりやすくなる．

\subsection{将来温度予測と不確実性の定量化}

カルマンフィルタにより得られた $\hat{\mathbf{x}}_{t|t}$ および $\mathbf{P}_{t|t}$ を初期値として，$\Delta$ 秒先の温度を逐次予測する．
ここで，$\Delta$ をサンプリング周期 $\Delta t$ で割ったステップ数を $m=\Delta/\Delta t$ とする．

$t$ 時点の推定結果から $m$ ステップ先を予測した状態を $\hat{\mathbf{x}}_{t+m|t}$，その誤差共分散を $\mathbf{P}_{t+m|t}$ とする．
予測温度は

\begin{equation}
\hat{T}_{t+m|t}
=
\left(
\hat{\mathbf{x}}_{t+m|t}
\right)_1
\end{equation}

で与えられる．
ここで，$\hat{T}_{t+m|t}$ は時刻 $t$ における状態推定結果に基づいて算出された $\Delta$ 秒先の温度予測である．

また，カルマンフィルタにより得られる誤差共分散行列 $\mathbf{P}_{t+m|t}$ は，将来予測の不確実性を表す指標として利用できる．
本研究では，温度成分に対応する分散

\begin{equation}
P_{t+m|t}^{(T)}
=
(\mathbf{P}_{t+m|t})_{1,1}
\end{equation}

を不確実性指標として用いる．

対応する標準偏差は以下のようになる。

\begin{equation}
\sigma_{t+m|t}^{(T)}
=
\sqrt{P_{t+m|t}^{(T)}}
\end{equation}


この値は，観測とモデルのずれが大きい場合や外部熱影響の推定が不安定な場合に増加するため，将来温度予測の信頼度低下を表す指標として解釈できる．
従来手法のように過去誤差の固定統計量に依存するのではなく，推定状態に応じた時変な不確実性評価が可能となる点が本手法の特徴である．
また，現在の不確実性はフィルタ後共分散 $\mathbf{P}_{t|t}$ により評価され，将来予測時にはこれを時間発展させた分散 $\mathbf{P}_{t+m|t}$ を用いる．

\subsection{動的マージン決定}

従来手法では固定マージンを用いるため，予測誤差の時間変動に十分対応できないという問題がある．
本研究では，上述の予測不確実性に応じて動的にマージンを決定する．

提案手法のマージンを次式で定義する．

\begin{equation}
M_t
=
\alpha \sigma_{t+m|t}^{(T)} + \delta
\end{equation}

ここで，$\alpha$ は安全係数，$\delta$ は実環境における不確実性の過小評価を補正するための定数項である．
$\alpha \sigma_{t+m|t}^{(T)}$ は理論的な不確実性に基づく成分であり，$\delta$ はモデル化誤差や未観測要因に起因する不確実性を補償するための実用的な補正項として導入する．
本研究では，評価実験において $\alpha=2.0$，$\delta=0.5$ を用いた．

このとき，安全温度閾値を以下のように定義する．

\begin{equation}
T_{\mathrm{safe},t}
=
T_{\mathrm{limit}} - M_t
\end{equation}

ここで，$T_{\mathrm{limit}}$ は物理的または運用上の上限温度である．
不確実性が大きい場合にはマージンが拡大し，温度逸脱リスクを低減する．
一方，不確実性が小さい場合にはマージンが縮小し，リソース利用効率の向上を図る．
このため，本手法は安全性を保ちながらも，不必要に厳しい制御を避けやすい．

なお，本手法は過去の予測誤差統計に基づいてマージンを調整する坂本型手法とは異なり，外部熱影響そのものを状態としてモデル化し，その推定不確実性を用いてマージンを決定する点に特徴がある．
すなわち，本手法は誤差を事後的に補正するのではなく，誤差の原因を状態として扱い，予測段階で考慮する点に本質的な違いがある．

\subsection{制御方針}

動的に決定された安全閾値 $T_{\mathrm{safe},t}$ に基づき，ノードの負荷制御を行う．
$t$ 時点で得られた $\Delta$ 秒先温度予測 $\hat{T}_{t+m|t}$ が安全閾値を超える場合，

\begin{equation}
\hat{T}_{t+m|t} > T_{\mathrm{safe},t}
\end{equation}

となるため，制御を発動する．

本手法では，将来温度の予測結果に基づき，
温度が上限に到達する前に負荷を抑制する点に特徴がある．
すなわち，実測温度が閾値を超過した後に制御を行うのではなく，
将来の温度上昇を見越して事前に制御を行うことで，
温度逸脱の抑制を図る．

\section{評価実験}

本章では，提案手法の有効性を評価するための実験条件および評価方法について述べる．
本研究では，（1）空間熱干渉の存在確認，（2）温度予測精度の評価，（3）動的マージン制御の有効性の評価，の3点について検証した．

\subsection{評価目的}

本研究の目的は，空間熱干渉によって生じる予測誤差を考慮し，温度逸脱を抑制しつつ過度なリソース制限を回避することである．
そのため，以下の観点から評価を行う．

\begin{enumerate}
    \item 隣接ノード負荷が対象ノード温度へ与える影響の確認
    \item 外部熱影響を考慮することによる温度予測精度の改善
    \item 動的マージンにより安全性と効率の両立が可能か
\end{enumerate}

\subsection{実験環境}

実験には Raspberry Pi 4B を2台使用し，密閉されたプラスチック筐体内に高密度に配置した．
ノード間距離は約5\,cm とした．
対象ノードを Raspi1，隣接ノードを Raspi2 とする．

各ノードのCPU温度およびCPU使用率は内蔵センサおよびOS情報から取得し，MQTTを用いて収集ノードへ送信した．
生データの取得周期は約1秒であるが，温度変化は本環境において時定数が約60秒程度の比較的緩やかな現象であり，
1秒周期のデータには観測ノイズや短周期の揺らぎが多く含まれる．
そこで，本研究ではノイズ低減および計算安定性の向上を目的として，
5秒区間ごとに平均化したデータを用いてモデル化および状態推定を行った．

アイドル状態とは，負荷生成プログラムを停止し，OSのバックグラウンド処理のみが動作している状態と定義する．
負荷生成には \texttt{stress-ng} を用いた．

\subsection{比較手法}

提案手法の有効性を評価するため，以下の3手法を比較対象とした．

\begin{itemize}
    \item Self-only モデルでは

\begin{equation}
T_{t+1} = aT_t + bu_t + d
\end{equation}

を用いる．
このような温度ダイナミクスは，熱伝導を電気回路に対応付けたRCモデルとして表現できることが知られており，温度は指数的な応答特性を持つシステムとして近似可能である~\cite{skadron2003temperature}．

    \item 坂本型手法  

    Self-only モデルによる予測値に対し，
    過去の予測誤差系列から算出した統計量に基づいてマージンを補正する手法である．
    外部熱影響は状態として扱わず，誤差として事後的に補正する．

    \item 提案手法 
    
    第3章で述べた状態空間モデルとカルマンフィルタを用い，
    外部熱影響を状態として推定するとともに，
    その不確実性に基づいて動的にマージンを決定する手法である．
\end{itemize}

\subsection{実験条件}

本研究では目的に応じて2種類の負荷パターンを用いた．

空間熱干渉の観測では，5分アイドル，10分高負荷，5分冷却の負荷パターンを用いた．
このとき，対象ノード Raspi1 はアイドル状態とし，隣接ノード Raspi2 のみに高負荷を印加した．
これにより，自ノード負荷がほぼない場合でも，隣接ノードの発熱により対象ノード温度が変化するかを観測した．

温度予測および制御評価では，5分アイドル，10分負荷，5分冷却の負荷パターンを用いた．
前半および後半では Raspi1 のみに負荷を印加し，中央の10分では Raspi1 と Raspi2 の両方に負荷を印加した．
これにより，外部熱影響の投入および解除を含む条件を再現した．

\subsection{モデルとパラメータ設定}

Self-only モデルでは

\begin{equation}
T_{t+1} = aT_t + bu_t + d
\end{equation}

を用いる．

提案手法では，第3章で示した拡張モデルを用い，
隣接ノード温度を中心化した入力

\begin{equation}
z_t = z_t^{\mathrm{nbr}} - z_{\mathrm{base}}
\end{equation}

を導入する．

係数 $a,b,c,d$ は観測データに対して最小二乗法を適用して求めた．
本評価で得られた主な係数は表\ref{tab:model_param}に示す通りである．また、本研究で用いたカルマンフィルタおよび適応ノイズの主要パラメータを表\ref{tab:kf_param}に示す．

\begin{table}[tbp]
\centering
\caption{推定されたモデル係数}
\label{tab:model_param}
\begin{tabular}{lcccc}
\hline
モデル & $a$ & $b$ & $c$ & $d$ \\
\hline
Self-only & 0.947 & -1.701 & -- & 5.181 \\
Neighbor-aware & 0.899 & -1.625 & 0.031 & 8.504 \\
\hline
\end{tabular}
\end{table}



\begin{table}[tb]
\centering
\caption{カルマンフィルタおよび適応ノイズのパラメータ}
\label{tab:kf_param}
\begin{tabular}{lc}
\hline
パラメータ & 値 \\
\hline
減衰係数 $\rho$ & 0.90 \\
温度プロセスノイズ $q_T$ & 0.038 \\
外部熱影響プロセスノイズ $q_H^{(0)}$ & 0.008 \\
観測ノイズ $R$ & 0.077 \\
適応係数 $\beta_{qh}$ & 2.0 \\
\hline
\end{tabular}
\end{table}

これらの値は，予測誤差および制御性能を比較しながら経験的に調整したものであり，
特に $\beta_{qh}$ は外乱変動に対する追従性と安定性のトレードオフを考慮して決定した．

また，本研究では予測ホライズンを $\Delta=30$ 秒と設定した．
対象ノードに負荷を印加した際の温度応答から，
本システムの時定数は約 66 秒程度であることが確認された．
一般に，時定数と同程度の時間スケールでは温度変化が顕著に現れるため，
将来温度の予測においても同程度のオーダーの先読みが有効である．
一方で，過度に長い予測は誤差の蓄積を招くため，
本研究ではその約半分に相当する 30 秒を予測時間として採用した．

さらに，上限温度は $T_{\mathrm{limit}}=80\,^\circ\mathrm{C}$ と設定した．
これは，Raspberry Pi 4B において CPU 温度が約$ 80\,^\circ\mathrm{C}$ に達すると
サーマルスロットリングが発生し，性能低下が生じるためである．
本研究では，この温度を運用上の制約温度として用いた．
\subsection{評価指標}

温度予測精度の評価には，以下の指標を用いた．

\begin{itemize}
    \item 平均絶対誤差（MAE）
    \item 二乗平均平方根誤差（RMSE）
    \item 決定係数（$R^2$）
\end{itemize}

また，制御評価では以下の指標を用いた．

\begin{itemize}
    \item 制御発動割合（control ratio）
    \item 平均マージンおよび最大マージン
    \item 温度逸脱回数（miss count）
\end{itemize}

ここで，温度逸脱は実測温度が上限温度 $T_{\mathrm{limit}}$ を超えた場合と定義する．
また，制御判断には第3章で定義した将来予測温度 $\hat{T}_{t+m|t}$ を用いた．

これらの指標により，温度予測精度と制御挙動の両面から手法の有効性を評価する．
温度予測精度に関する指標は，将来温度の推定誤差の大きさおよび傾向を定量的に評価することを目的としている．
一方，制御評価指標は，温度逸脱の抑制とリソース利用効率のトレードオフを評価するために用いる．


\begin{figure}[tb]
    \centering
    \includegraphics[width=1.0\linewidth]{image-7.png}
    \caption{隣接ノード負荷による温度変化}
    \label{fig:temp_change}
\end{figure}
\section{評価結果}

\subsection{隣接ノード負荷による空間熱干渉の観測}

隣接ノードに負荷を与えた際の温度変化を図\ref{fig:temp_change}に示す．



図\ref{fig:temp_change}より，隣接ノードに負荷を印加した後，
対象ノードは自身に負荷を与えていないにもかかわらず温度が徐々に上昇していることが確認できる．
この温度上昇は即時的ではなく，時間遅れを伴って緩やかに進行しており，
熱の蓄積および伝播による動的な影響であると解釈できる．
この結果は，隣接ノードの発熱が自ノード温度に影響を与える空間熱干渉が実際に存在することを示している．
したがって，自ノード負荷のみに基づく従来モデルでは，
このような外部起因の温度変動を十分に説明できない可能性がある．

\subsection{温度予測精度の評価}

30秒先温度予測における各手法の精度比較結果を表\ref{tab:pred_result}に示す．
\begin{table}[tbp]
\centering
\caption{30秒先温度予測精度}
\label{tab:pred_result}
\begin{tabular}{lrrrr}
\hline
手法 & Bias [℃] & MAE [℃] & RMSE [℃] & $R^2$ \\
\hline
Self-only & 0.041 & 0.999 & 1.290 & 0.913 \\
坂本型 & -0.074 & 1.164 & 1.489 & 0.884 \\
提案手法 & 0.131 & 0.878 & 1.139 & 0.932 \\
\hline
\end{tabular}
\end{table}
各手法の予測温度と実測温度の比較を図\ref{fig:pred_compare}に示す．
図の可読性向上のため，観測温度および各予測温度には6点の移動平均を適用した．
図からも，提案手法が全体として実測温度に良好に追従している傾向が確認でき，この傾向は表\ref{tab:pred_result}の定量評価結果と整合している．
\begin{figure}[tb]

    \centering

    \includegraphics[width=1.00\linewidth]{image-9.png}

    \caption{実測温度と各手法の30秒先予測温度の比較（6点移動平均）}

    \label{fig:pred_compare}

\end{figure}


表\ref{tab:pred_result}より，提案手法はSelf-onlyモデルおよび坂本型手法と比較して，
MAEおよびRMSEの両方で最も小さい値を示している．また，$R^2$も最も高い値となっており，
全体として予測性能が向上していることが確認できる．

この結果は，隣接ノード由来の熱影響を外部熱状態として明示的にモデル化し，
カルマンフィルタにより逐次推定を行うことで，
温度変動の背後にある要因を適切に捉えられたためであると考えられる．
一方，坂本型手法は過去の誤差統計に基づく補正であるため，
外部要因の変動を直接的に扱うことができず，
本実験のような環境では精度が劣化したと考えられる．

\subsection{動的マージン制御の評価}

各手法におけるマージンおよび制御指標の比較結果を表~\ref{tab:margin_result}に示す．
さらに，各手法におけるマージンの時間変化を図~\ref{fig:margin}に示す．
\begin{table}[tbp]
\centering
\caption{マージンおよび制御指標の比較}
\label{tab:margin_result}
\begin{tabular}{lr@{ }rrr}
\hline
手法 & \shortstack{平均マージン \\ {[℃]}} & \shortstack{最大マージン \\ {[℃]}} & 制御率 & \shortstack{温度逸脱 \\ 回数} \\
\hline
Fixed & 1.861 & 1.861 & 0.809 & 0 \\
坂本型 & 1.827 & 4.062 & 0.763 & 0 \\
提案手法 & 1.491 & 1.823 & 0.750 & 0 \\
\hline
\end{tabular}
\end{table}

\begin{figure}[tb]

    \centering

    \includegraphics[width=1.0\linewidth]{image-8.png}

    \caption{各手法におけるマージンの時間変化}

    \label{fig:margin}

\end{figure}
表~\ref{tab:margin_result}および図~\ref{fig:margin}より，
固定マージン手法では一定の安全余裕を維持する一方で，
状況に応じた調整が行えず，常に過大なマージンとなる．
坂本型手法では平均マージンは低減されるものの，
外乱発生時にはマージンが大きく変動し，
最大で4℃を超える過大な値となる区間が確認できる．

これに対し提案手法では，
マージンは比較的安定して推移しつつ，
必要な区間においてのみ適度に増加している．
これは，カルマンフィルタによる不確実性評価に基づき，
時刻ごとにマージンを調整しているためである．


\section{考察}

本研究では，空間熱干渉を考慮した状態推定に基づく動的熱マージン制御手法を提案し，
その有効性を評価した．本節では，得られた結果に基づき，
提案手法の特徴および従来手法との違いについて考察する．

まず，温度予測精度の観点では，
提案手法はSelf-onlyモデルおよび坂本型手法と比較して最も高い精度を示した．
この要因は，外部熱影響を状態変数としてモデルに組み込み，
カルマンフィルタにより逐次推定を行った点にある．
従来手法では外部要因は予測誤差として事後的に扱われるのに対し，
提案手法ではその原因を状態として取り込むことで，
予測段階から外乱の影響を反映できる構造となっている．
この違いにより，外部環境が変動する状況においても，
安定した予測性能が得られたと考えられる．

次に，マージン制御の観点では，
提案手法は平均マージンを抑えつつ温度逸脱を防止できることが確認された．
固定マージン手法は安全性を確保しやすい一方で，
常に過大な余裕を持つためリソース利用効率が低下する．
坂本型手法は平均マージンを低減できるが，
過去の誤差統計に依存するため，
環境変化に対して過剰または過小に反応する場合がある．

これに対し提案手法では，
カルマンフィルタの誤差共分散を用いて
時刻ごとの不確実性を直接評価し，
それに応じてマージンを調整する．
すなわち，坂本型手法が「過去の誤差」に基づくのに対し，
提案手法は「現在の不確実性」に基づいて制御を行う点に本質的な違いがある．

一方で，本実験環境では温度変動が比較的滑らかであり，
不確実性が大きく変動する区間が限定的であったため，
動的マージンの効果は平均マージンの低減という形で現れたものの，
固定マージン手法との差は限定的であった．
また，動的マージンの性能は不確実性推定の精度に依存するため，
モデル化誤差やノイズ設定の影響を受けやすい．
このため，本研究では補正項 $\delta$ を導入することで，
理論モデルと実環境の差異を吸収し，
実用的な安全性を確保した．

以上より，本研究の本質的な貢献は，
外部熱影響を状態として推定し，
その不確実性を制御判断に直接反映した点にあるといえる．
この枠組みにより，
環境変動を伴うシステムにおいても，
安全性と効率を両立した熱制御が可能となる．


\section{おわりに}

本稿では，エッジマイクロクラスタにおける空間熱干渉に着目し，
状態推定に基づく動的熱マージン制御手法を提案した．

提案手法では，隣接ノードからの熱影響を時間的に変動する状態量としてモデル化し，
カルマンフィルタにより自ノード温度と外部熱影響を同時に推定した．
さらに，推定誤差共分散に基づいて予測の不確実性を定量化し，
それに応じて安全マージンを動的に調整する枠組みを構築した．

評価実験の結果，
空間熱干渉が実際に温度へ影響を与えることを確認するとともに，
提案手法が従来手法と比較して高い温度予測精度を示すことを確認した．
また，動的マージンにより，
温度逸脱を防止しつつ平均マージンを低減できることを示し，
安全性とリソース利用効率の両立が可能であることを示した．

一方で，動的マージンの効果は環境条件に依存するため，
より急激な外乱や複雑な熱干渉が存在する状況における評価が今後の課題である．
また，パラメータ同定の自動化や，
DVFSおよびタスクマイグレーションとの統合による実システムへの適用も重要な課題である．

本手法は，外乱が時間的に変動するシステムにおいて，
不確実性を考慮した制御判断を可能とする一般的な枠組みであり，
実環境における温度変動に対して，
より柔軟に対応可能な熱管理手法として有効であると考えられる．

\bibliographystyle{ipsjunsrt}
\bibliography{references}

\end{document}